% =========================================================================
% General Introduction
% =========================================================================
\chapter*{General Introduction}
\addcontentsline{toc}{chapter}{General Introduction}
\markboth{General Introduction}{}

Short-form video --- TikTok clips, Instagram Reels and YouTube Shorts under sixty seconds ---
has become the dominant content format of the social-media era. According to Sensor Tower's
2024 \textit{State of Mobile} report, TikTok crossed 1.5~billion monthly active users and
short-form video formats now generate the majority of advertising revenue on every major
mobile platform. For an independent creator or a small marketing team, the economic
opportunity is real, but so is the cost of guessing wrong: the median Short receives fewer than
1{,}000 views, while a viral hit can reach millions, and the difference is decided in the first
three seconds by a combination of visual hook, audio rhythm, on-screen captions and pixel
quality.

Despite that asymmetry, almost no tooling exists today that helps a creator answer the most
important questions \emph{before} they hit publish:

\begin{enumerate}
  \item Will this clip perform? Which 473-feature combination of visuals, audio, transcript
        and tone correlates with virality, and which one is the bottleneck on \emph{this}
        particular video?
  \item Are the burned-in captions readable on a 6-inch phone, or are the long words clipping
        the safe area?
  \item Is the footage --- often shot on a hand-held smartphone, in low light, with motion blur ---
        being uploaded at the visual quality the algorithm and the audience expect?
\end{enumerate}

Existing post-publication analytics tools (TubeBuddy, VidIQ) and post-production editors
(CapCut, Descript) address one of these questions in isolation; none of them brings the three
together with an integrated multi-platform publishing layer.

This end-of-studies project, carried out as part of the Bachelor / Engineering programme at
\textit{[Institution]}, conceives and builds \textbf{BViral Control Center}: a four-pillar AI
platform that closes that gap. The system fuses pretrained vision and language models, classical
computer vision, gradient-boosted trees and a typed full-stack web application into a single
workflow that ingests a raw clip and returns a virality score with explanations, a captioned
version with adaptive font sizing, an enhanced version with super-resolution and face
restoration, and a one-click cross-platform scheduler.

\bigskip

To attain that goal, this internship report is divided into three major parts:

\begin{itemize}
  \item \textbf{Chapter 1 --- State of the Art:} this chapter focuses on the analysis of the
        existing solutions for short-form video creators, the functional and non-functional
        requirements of the proposed platform, the major actors involved, and the working
        methodology adopted for the project.
  \item \textbf{Chapter 2 --- Conceptual Study:} a detailed conception is given through UML
        diagrams (use cases, class diagram, sequence diagrams) and the dataset and AI pipeline
        specifications, a major step before implementing the solution.
  \item \textbf{Chapter 3 --- Realisation:} this chapter justifies the choice of technologies,
        documents the implementation of each of the four AI pillars together with concrete
        evaluation metrics and code listings, and presents screenshots of the application.
\end{itemize}

Finally, we conclude with a general conclusion in which we discuss the future perspectives of
the platform and the potential changes intended to further improve the system and make it more
useful for short-form video creators of every size.

\cleardoublepage
